\section{Rank-two systems and integral congruences}\label{sec:systems}

The cyclic covers now supply the Galois representations through which we
will transfer automorphy.  The involution identifies opposite rotation
characters, giving each character space the determinant and Hodge--Tate
weights of an elliptic curve.  The integral regular-module calculation
will then give congruences even at primes dividing the covering degree.

Return to an imaginary quadratic field $K$ and a non-CM elliptic curve
$E/K$.  By Serre's open image theorem \cite[\S4.4, Th\'eor\`eme~3]{SerreOpenImage}, choose a prime
$p>100$ such that the image of $G_K$ on $E[p]$ contains
$\SL_2(\F_p)$.  Set
\begin{equation}\label{sys:standing-fields}
 n=3p,\qquad F=K(\mu_n).
\end{equation}
Both $K/\Q$ and $\Q(\mu_n)/\Q$ are abelian, so $F/\Q$ is abelian.
It is a CM field: complex conjugation is central, its fixed field is
totally real, and $F$ has no real embedding because it contains $K$.
Moreover,
\begin{equation}\label{sys:no-fifth-roots}
 \zeta_5\notin F.
\end{equation}
Indeed, $\Q(\mu_n)$ is unramified at $5$, and the ramification index
of $F$ at $5$ is consequently at most $[K:\Q]=2$.  A field containing
$\Q(\zeta_5)$ has ramification index divisible by $4$ at $5$.

Apply the construction of Lemma~\ref{geo:family} with this $n$.
Write $U\subset\A^1_K$ for its parameter open and $B_t$ for the
quadratic twist of $E$ used there.  We may shrink $U$ whenever a smooth
proper model of one of the finitely many families is required.  Until
Section~\ref{sec:residual}, the point $t\in U(K)$ is arbitrary.

\subsection{Character spaces}

Fix an embedding $\overline\Q\hookrightarrow\C$ and put
$M=\Q(\mu_n)\subset\overline\Q$.  For each rational prime $\ell$, choose
an embedding
$j_\ell:\overline\Q\hookrightarrow\overline\Q_\ell$ and an isomorphism
$\iota_\ell:\overline\Q_\ell\xrightarrow{\sim}\C$ whose composition
with $j_\ell$ is the fixed complex embedding.  Every rotation character
in this section is a homomorphism
\[
 \theta:\mu_n(F)\longrightarrow M^\times.
\]
Write $C_m(t)=C_{m,t}$ for a fiber of the family.  For the curve
$C_n=C_n(t)$, define
\begin{equation}\label{sys:definition}
 V_{\theta,\ell}
 =\bigl\{x\in H^1_{\mathrm{\acute et}}
       (C_{n,\overline F},\overline\Q_\ell):
       u^*x=j_\ell(\theta(u))x\text{ for all }u\in\mu_n(F)\bigr\}.
\end{equation}
Here $u$ denotes the deck transformation with multiplier $u$.
All deck transformations are defined over $F$, so these are
$G_F$-representations.  We write $V_\theta$ for the collection of these
realizations.  In a congruence at $\ell$, all characters use the same
embedding $j_\ell$; this convention is essential when comparing their
reductions.

\begin{proposition}\label{sys:properties}
For every $t\in U(K)$ and every rotation character $\theta$, the
representations $V_{\theta,\ell}$ have the following properties.
\begin{enumerate}
\item They are semisimple of dimension two, with determinant
$\eps_\ell^{-1}$.
\item At every place of $F$ above $\ell$, they are potentially
semistable with labelled Hodge--Tate weights $\{0,1\}$.
\item There is a finite set of places $S$ of $F$, independent of
$\ell$, such that $V_{\theta,\ell}$ is unramified at every
$v\notin S$ of residue characteristic different from $\ell$.
At such a place its Frobenius characteristic polynomial is the image
under $j_\ell$ of a polynomial in $M[T]$ independent of $\ell$.
\end{enumerate}
\end{proposition}

\begin{proof}
The rank-two regular-module statement of
Proposition~\ref{mon:integral-module}, followed by comparison between
singular and \(\ell\)-adic cohomology, gives
$\dim V_{\theta,\ell}=2$.

Let $\langle\ ,\ \rangle$ be the cup-product pairing on $H^1(C_n)$,
with values in $\overline\Q_\ell(-1)$.  It pairs the $\theta$-space
perfectly with the $\theta^{-1}$-space.  The involution $w$ is defined
over $F$, preserves cup product, and exchanges these two spaces.
Consequently
\[
 [x,y]_\theta=\langle x,w^*y\rangle,
 \qquad x,y\in V_{\theta,\ell},
\]
is a perfect $G_F$-equivariant pairing into
$\overline\Q_\ell(-1)$.  It is alternating, because
\[
 \langle y,w^*x\rangle
   =\langle w^*y,x\rangle
   =-\langle x,w^*y\rangle.
\]
A perfect alternating pairing on a two-dimensional space identifies
its determinant with the target line.  This proves
$\det V_{\theta,\ell}=\eps_\ell^{-1}$.

To compute the Hodge numbers, make any embedding of $F$ into $\C$.
Let $a_\theta$ be the multiplicity of $\theta$ in the holomorphic
differentials of the resulting complex curve.  The holomorphic
involution $w$ gives $a_\theta=a_{\theta^{-1}}$.  Hodge decomposition
and complex conjugation identify the dimension of the full
$\theta$-space with $a_\theta+a_{\theta^{-1}}$.  Since this dimension
is two, $a_\theta=1$.  This argument works for every embedding and
every character.  The comparison theorems therefore give the labelled
weights $0,1$.  Potential semistability follows from semistable
reduction and the semistable comparison theorem for smooth proper
varieties; see \cite{Tsuji}.  These properties pass to the character
summands.

Outside a finite set, $C_n$ has good reduction and the rotations
extend to its smooth proper model.  Smooth proper base change gives
unramifiedness.  The projector onto the $\theta$-space is
\[
 e_\theta=\frac1n\sum_{u\in\mu_n(F)}
              j_\ell(\theta(u))^{-1}u^*.
\]
This formula is used only in characteristic zero.  The Lefschetz
trace formula for a rotation composed with a positive power of
Frobenius shows that its trace on $H^1(C_n)$ is independent of
$\ell$.  Applying the displayed projector gives traces in $M$,
compatible under the embeddings $j_\ell$.  Together with the
already determined determinant, this gives the asserted common
characteristic polynomial.

Finally, $H^1(C_n)$ is the first cohomology of its Jacobian.  Faltings'
semisimplicity theorem \cite[\S5, Satz~3]{Faltings} makes this a semisimple
$G_F$-representation.  Extension of scalars and passage to a
$G_F$-stable direct summand preserve semisimplicity, proving the
remaining assertion.
\end{proof}

\subsection{The two elliptic-curve realizations}

The characters of orders one and three will connect the new systems
to elliptic curves.  For the order-three character, the involution
produces a genus-one quotient rather than merely an abstract
rank-two summand.

\begin{proposition}\label{sys:elliptic}
If $\theta$ has order $m>1$ dividing $n$, the system $V_\theta$ is
realized by the corresponding character space in $H^1(C_m)$.
For the trivial character,
\[
 V_{1,\ell}\simeq H^1_{\mathrm{\acute et}}
                  (B_{t,\overline F},\overline\Q_\ell).
\]
Let $D_t=C_3(t)/\langle w\rangle$ and $A_t=\Jac(D_t)$.  Then $D_t$
has genus one, $A_t$ is an elliptic curve over $F$, and for either
nontrivial order-three character $\eta$ there are $G_F$-equivariant
isomorphisms
\[
 V_{\eta,\ell}\simeq
 H^1_{\mathrm{\acute et}}(A_{t,\overline F},\overline\Q_\ell)
 \qquad\text{for every }\ell.
\]
\end{proposition}

\begin{proof}
The quotient map $C_n\to C_m$ sends the covering coordinate $Y$ to
$Y^{n/m}$.  Pullback is injective on rational cohomology, since its
composition with trace is multiplication by the degree.  Its image
is the subspace fixed by the kernel of $\mu_n\to\mu_m$.
The rank-two calculation therefore identifies its order-$m$
character summand with $V_{\theta,\ell}$.  The same argument for the
quotient $C_n\to B_t$ identifies the trivial character space.

On $H^1(C_3)$, the trivial rotation space is $H^1(B_t)$, and $w$
acts there by $[-1]^*=-1$.  It exchanges the two nontrivial rotation
spaces, each of dimension two.  Their direct sum consequently has
a two-dimensional $w$-invariant subspace.  Rational cohomology of a
finite quotient is the invariant subspace, so
\[
 H^1(D_t)\simeq H^1(C_3)^{w=1}
\]
has dimension two.  Thus $D_t$ has genus one and its Jacobian is an
elliptic curve, whether or not $D_t$ has an $F$-rational point.
Projection of a vector $x+w^*x$ onto the $\eta$-space is $x$;
therefore this projection is a $G_F$-equivariant isomorphism from
the invariant subspace to $V_{\eta,\ell}$.  Finally,
$H^1(D_t)\simeq H^1(\Jac(D_t))$.  The latter identification is
independent of the geometric base point used in an Abel--Jacobi
map, because translations act trivially on cohomology.
\end{proof}

\subsection{Evaluation of the integral group ring}

A character projector contains the denominator $n$ and cannot by
itself establish a congruence at $3$ or $p$.  Instead, we use the
integral freeness proved in Section~\ref{sec:monodromy}.  Evaluation
of a free group-ring module gives a lattice before reduction; two
characters with the same residual values then give the same
residual representation.

\begin{lemma}\label{sys:integral-evaluation}
Fix $t\in U(K)$ and a prime $\ell$.  Let $L/\Q_\ell$ be a finite
extension containing $j_\ell(M)$, with ring of integers
$\mathcal O$, maximal ideal $\mathfrak m$, and residue field $k$.
For a rotation character $\theta$, the module
\[
 \mathcal T_\theta=
 H^1_{\mathrm{\acute et}}(C_{n,\overline F},\Z_\ell)
 \otimes_{\Z_\ell[\mu_n],\,j_\ell\theta}\mathcal O
\]
is a free rank-two $G_F$-stable lattice realizing
$V_{\theta,\ell}$ after extension to $\overline\Q_\ell$.
If $j_\ell\theta$ and $j_\ell\theta'$ have the same reduction as
characters $\mu_n\to k^\times$, then
\[
 \mathcal T_\theta/\mathfrak m\mathcal T_\theta
 \simeq
 \mathcal T_{\theta'}/\mathfrak m\mathcal T_{\theta'}
\]
as $k[G_F]$-modules.
\end{lemma}

\begin{proof}
By Proposition~\ref{mon:integral-module} and comparison,
\[
 H^1_{\mathrm{\acute et}}(C_{n,\overline F},\Z_\ell)
 \simeq \Z_\ell[\mu_n]^2
\]
as group-ring modules.  The Galois action commutes with this group
ring because the rotations are defined over $F$.  Tensoring by
evaluation at $j_\ell\theta$ therefore gives a free rank-two
$\mathcal O$-module with continuous Galois action.  In characteristic
zero the group algebra is semisimple; its quotient by character
evaluation is precisely the corresponding character summand.
This identifies its rational representation with $V_{\theta,\ell}$.

After reduction, the displayed tensor product becomes
\[
 H^1_{\mathrm{\acute et}}(C_{n,\overline F},\Z_\ell)
 \otimes_{\Z_\ell[\mu_n],\,\overline\theta} k.
\]
It depends only on the residual evaluation map
$\overline\theta$.  This proves the asserted isomorphism without
assuming that $\ell$ is prime to $n$.
\end{proof}

Choose a primitive character $\chi$ of order $n=3p$ and write its
unique primary decomposition as
\[
 \chi=\psi\eta,\qquad \operatorname{ord}(\psi)=p,
 \qquad \operatorname{ord}(\eta)=3.
\]
For a system $V_\theta$, write $\overline V_{\theta,\ell}$ for the
semisimplified reduction of a stable lattice, with scalars extended
to $\overline\F_\ell$.  This representation is independent of the
choice of stable lattice.

\begin{corollary}\label{sys:congruences}
For every $t\in U(K)$ there are isomorphisms
\begin{equation}\label{sys:congruence-chain}
 \overline V_{\eta,p}\simeq\overline V_{\chi,p},
 \qquad
 \overline V_{\chi,3}\simeq\overline V_{\psi,3},
 \qquad
 \overline V_{\psi,p}\simeq\overline V_{1,p}.
\end{equation}
These hold with the elliptic-curve realizations of
Proposition~\ref{sys:elliptic} at the two ends.
\end{corollary}

\begin{proof}
A root of unity of $\ell$-power order reduces to $1$ in residue
characteristic $\ell$.  Thus $\chi$ and $\eta$ have the same
reduction at $p$, $\chi$ and $\psi$ have the same reduction at $3$,
and $\psi$ and $1$ have the same reduction at $p$.
Lemma~\ref{sys:integral-evaluation} gives the three congruences.
The rational isomorphisms with the smaller covers and elliptic
curves in Proposition~\ref{sys:elliptic} preserve semisimplified
reduction, even if their natural integral pullbacks have index
divisible by $\ell$.
\end{proof}

The direction of propagation is now fixed: automorphy will pass
from $V_\eta$ to $V_\chi$ at $p$, then to $V_\psi$ at $3$, and
finally to $V_1$ at $p$.  The congruences themselves impose no
restriction on $t$.  The subsequent local and residual arguments will
choose one parameter for which all three automorphy-lifting steps apply.
